\scanpage{244}
Si on ne fait pas d'hypothèse supplémentaire, il n'est pas sûr que l'opération de $Z$ sur $\pi$ soit fidèle, i.e. qu'on ait

\[
\left\{\begin{aligned}
&\forall\lambda\in Z\ \deleted{\text{tel qu'il existe}}\ u\in H\ \deleted{\text{avec}}\\
&(s_0,\lambda t_0)=u(s_0,t_0)\qquad\Longrightarrow\lambda=1\quad.\\
&\text{i.e. }\quad us_0=s_0,\ u(t_0)=\lambda t_0
\end{aligned}\right.
\tag*{(72)}
\]

Pour ceci, il suffit que tout sommet soit déterminé par son étoile --- condition qui ne s'exprime pas, il me semble, de façon particulièrement jolie en termes des données $\unsure{S|Z};H\subset\unsure{\mathfrak S|\mathcal G}_S,\ (s_0,t_0)\in S\times S$.

Je vais, pour simplifier, supposer qu'on a (72), --- sinon dans (64), l'hom $\unsure{\delta_1|\delta,}$ va de \deleted{\unsure{$H$}} $G$, non dans $Z$, mais dans le groupe quotient $Z/\unsure{\mathfrak z|\zeta}_Z$, où $\mathfrak z_Z$ est le noyau de l'opération de $Z$ sur $\pi$, i.e le sous-groupe formé des $\lambda\in Z$ pour lesquels $\exists u\in H$ avec $u(s_0,t_0)=(s_0,\lambda t_0)$.

Posons

\[
\left\{\begin{aligned}
U_0&=\deleted{\unsure{G|O}}H_{s_0}=\{u\in H\mid us_0=s_0\}\\
U_1&=H_{t_0}=\{u\in H\mid ut_0=\unsure{t|t_0}\}\\
\mathfrak z_H&=U_0\cap U_1=H_{(s_0,t_0)}=\{u\in H\mid us_0=s_0,ut_0=t_0\,,
\end{aligned}\right.
\tag*{(73)}
\]

et soit

\[
\sigma\in H\qquad\sigma(s_0,t_0)=(t_0,s_0)\quad\text{i.e.}\quad\sigma s_0=t_0,\sigma t_0=s_0.
\tag*{(74)}
\]

On a donc

\[
\left\{\begin{aligned}
U_1&=\sigma(U_0),\quad U_0=\sigma(\deleted{\unsure{H}}U_1)\\
\sigma^2&\in\mathfrak z_H\overset{\mathrm{def}}{=}U_0\cap U_1\ \deleted{\unsure{\sim}}\quad.
\end{aligned}\right.
\tag*{(75)}
\]

% archive sheet number: 243
